\section{不同指令怎样越过各自的提交边界}

“处理器执行一条指令”不能只画成一个方框。算术指令、装入、存储、分支和返回依赖的外部
状态不同，失败时允许发生的效果也不同。教学处理器按指令类别规定提交规则。

\begin{longtable}{|L{0.18\linewidth}|L{0.25\linewidth}|L{0.29\linewidth}|L{0.18\linewidth}|}
\hline
\rowcolor{BookTableHead}
指令类别 & 完成前读取 & 候选效果 & 提交条件 \\ \hline
寄存器算术 & 源寄存器、立即数 & 目标寄存器与标志的新值 & 当前指令无译码错误 \\ \hline
装入 & 基址寄存器、RAM 响应 & 目标寄存器的新值 & 读响应为 \code{OK} \\ \hline
存储 & 基址、源寄存器 & RAM 中若干字节的新值 & 写响应为 \code{OK} \\ \hline
条件分支 & 测试寄存器、位移 & 顺序 PC 或目标 PC & 目标满足对齐约束 \\ \hline
\code{RET} & SP、栈读响应 & 新 PC 与增加后的 SP & 栈读成功且返回 PC 合法 \\ \hline
\end{longtable}

算术指令的结果在写回前只是处理器内部候选值。装入多了一笔外部读，不能在响应之前写目标
寄存器。存储的效果发生在 RAM 一侧；处理器收到成功响应后才能继续，但不能先改 RAM、
随后又向软件报告该指令未执行。教学模型要求目标对每笔写要么完整接受指定字节，要么返回
错误且不改变任何字节。

\subsection{装入等待期间保存什么}

假设 RAM 需要三个周期完成读取。处理器至少要保留下列临时事实：

\begin{lstlisting}
pending_load:
  instruction_pc = 0x00100004
  destination    = r5
  address        = 0x00102000
  width          = 4
  next_pc        = 0x00100008
\end{lstlisting}

它们不是软件可直接读取的通用寄存器，而是处理器内部执行状态。若等待时把 PC 提前提交到
\code{0x00100008}，又允许下一条指令读取旧 \code{r5}，程序会观察到错误顺序。最小处理器
因此在响应返回前不开始下一条指令。

\begin{center}
\begin{tikzpicture}[x=1.55cm,y=0.72cm]
\foreach \x in {0,1,2,3,4,5} {
  \draw[dashed,BookRule] (\x,0) -- (\x,3.8);
  \node[font=\scriptsize,BookMuted] at (\x,-0.25) {\x};
}
\node[anchor=east,font=\footnotesize] at (-0.25,3.2) {请求};
\node[anchor=east,font=\footnotesize] at (-0.25,2.2) {RAM};
\node[anchor=east,font=\footnotesize] at (-0.25,1.2) {提交};
\draw[very thick,BookAccent] (0,3.0) -- (1,3.0) -- (1,3.4) -- (2,3.4) -- (2,3.0) -- (5,3.0);
\node[font=\scriptsize] at (1.5,3.62) {接受};
\draw[very thick,BookMuted] (0,2.0) -- (2,2.0);
\draw[very thick,BookAccent] (2,2.0) -- (4,2.0);
\node[font=\scriptsize] at (3,2.28) {读取字节};
\draw[very thick,BookMuted] (4,2.0) -- (5,2.0);
\draw[very thick,BookMuted] (0,1.0) -- (4,1.0);
\draw[very thick,BookAccent] (4,1.0) -- (4,1.4) -- (5,1.4);
\node[font=\scriptsize,anchor=west] at (4.08,1.62) {\code{r5=7, PC+=4}};
\end{tikzpicture}
\end{center}
\diagramnote{请求被接受不等于装入已提交。RAM 完成读取后，处理器才同时更新目标寄存器和
顺序 PC。}

\subsection{存储的提交点位于目标响应}

\code{STORE r4,[r2+0]} 的处理器请求包含地址、宽度和数据。RAM 检查地址及字节使能，
在接受写入的边界更新四个字节并返回 \code{OK}。若地址未对齐，它返回
\code{ALIGNMENT} 且四个字节全部保持旧值。

对软件而言，存储指令完成后，同一处理器随后读取该地址必须看到新值。当前机器只有一个
请求发起者，没有缓存和其他处理器，所以不需要讨论多核可见顺序。这个简化条件应明确写出，
不能把单核结论直接推广到并发机器。

\subsection{分支先计算两个候选，再提交一个}

\code{BNEZ} 的顺序后继为 \code{PC+4}，目标为
\[
\texttt{PC+4}+4\times\operatorname{signextend}(\texttt{disp16}).
\]
处理器读取测试寄存器，在两个候选中选择一个写入 PC。未被选择的地址不会产生取指事务。
对循环分支，\code{disp16=0xfffb} 表示 \(-5\)，目标计算为
\[
\texttt{0x00100018}-20=\texttt{0x00100004}.
\]

若位移计算溢出 32 位地址或目标未对齐，教学处理器不提交普通 PC，而进入固定错误入口。
真实体系结构可能规定模地址运算或不同异常行为；本书只依赖“错误不会静默成为另一条合法
顺序指令”这一教学契约。

\subsection{\code{CALL} 与 \code{RET} 为什么同时改变内存和 PC}

\code{CALL target} 需要保留顺序后继，以便子程序结束后继续。教学处理器按以下顺序定义
调用：

\begin{enumerate}
\item 计算返回地址 \code{PC+4}；
\item 计算新栈指针 \code{SP-4}，检查是否对齐；
\item 向 \code{MEM32[SP-4]} 写返回地址；
\item 写响应成功后提交 \code{SP<-SP-4} 和 \code{PC<-target}。
\end{enumerate}

若第 3 步失败，SP 和 PC 都保持调用前的体系结构值。\code{RET} 进行相反方向的动作：读取
\code{MEM32[SP]}，成功后提交 \code{PC<-value} 和 \code{SP<-SP+4}。调用与返回因此不是
纯粹的 PC 跳转，它们还维护一条位于 RAM 的控制流历史。

\begin{longtable}{|L{0.20\linewidth}|L{0.29\linewidth}|L{0.29\linewidth}|}
\hline
\rowcolor{BookTableHead}
状态 & 调用前 & \code{CALL} 成功后 \\ \hline
\code{PC} & \code{0x00100100} & 子程序入口 \code{0x00100200} \\ \hline
\code{SP} & \code{0x001efffc} & \code{0x001efff8} \\ \hline
栈字 \code{0x001efff8} & 不属于有效栈 & 返回地址 \code{0x00100104} \\ \hline
\end{longtable}

当前求和核心没有调用子程序，但启动协议仍预置一个返回地址，使最外层 \code{RET} 能把
控制权交回固件。最外层返回与普通函数返回使用相同处理器规则，区别只在栈中字节由谁建立。

\subsection{错误响应怎样到达固定固件入口}

为了让错误至少可诊断，教学处理器预留三个只在错误入口使用的状态：

\begin{lstlisting}
fault_pc       address of the instruction that could not complete
fault_address  external address, if the instruction issued one
fault_cause    INVALID_OPCODE, UNMAPPED, ALIGNMENT, or READ_ONLY
\end{lstlisting}

错误发生时，处理器写入这些状态，把 PC 改为启动存储中的
\code{0x00000400}，并使用固件错误栈。失败指令的普通目标寄存器、普通顺序 PC 和目标内存
效果均不提交。固件错误入口读取三个字段，通过串行接口报告，然后复位机器。

\begin{center}
\begin{tikzpicture}[node distance=10mm and 15mm]
\node[stage,minimum width=3.1cm] (inst) {失败指令\\旧体系结构状态};
\node[stage,minimum width=3.3cm,right=of inst] (fault) {保存故障事实\\PC、地址、原因};
\node[stage,minimum width=3.3cm,right=of fault] (entry) {固定固件入口\\报告并复位};
\node[stage,minimum width=3.3cm,below=13mm of fault,fill=BookWarm] (no) {不提交\\普通写回与顺序 PC};
\draw[flow] (inst) -- (fault);
\draw[flow] (fault) -- (entry);
\draw[->,>=Latex,thick,draw=BookMuted] (fault) -- (no);
\end{tikzpicture}
\end{center}
\diagramnote{错误入口保存的是足以说明失败位置的最小事实，不是可稍后恢复的完整任务现场。
报告结束后复位，因此当前机制仍不支持任务级恢复。}

\subsection{无效操作码与未映射地址属于不同层次}

若取回的第一个指令字节不是已定义操作码，错误在处理器译码阶段产生。此时
\code{fault\_cause} 取值为 \code{INVALID\_OPCODE}，没有数据地址。若合法
\code{LOAD} 形成的地址没有
目标，处理器已经完成译码，互连返回 \code{UNMAPPED}，此时
\code{fault\_address} 有效。区分二者可以判断是代码字节本身无法解释，还是合法指令访问
了错误位置。

\subsection{为什么本章不把错误称为“异常处理系统”}

硬件已经能改变 PC 到固定入口，但入口只服务整机报告与复位。它没有独立任务身份，没有
保存全部通用寄存器，也没有返回失败指令之后继续执行的协议。后续建立受控入口时会复用
“保存原因并改变控制流”这一硬件能力，但不能提前把当前固定错误入口等同于完整操作系统
异常处理。
